La sonda solar \textit{Parker} (en anglès, \textit{Parker Solar Probe}) és una nau espacial en òrbita al voltant del Sol que té com a objectiu acostar-se molt a la superfície solar. La gràfica següent mostra com varia la distància de la nau respecte al Sol al llarg dels primers 1\,000 dies de missió i indica els instants A, B i C. Les unitats emprades per a mesurar la distància a la superfície del Sol són radis solars, $R_\mathrm{S}$.

\figura[0.75]{fig1.png}
{\centering\small\textsc{Font:} \textit{http://parkersolarprobe.jhuapl.edu}.\par}

\begin{apartat}{1,25}
Observeu a la gràfica els moments de màxim acostament al Sol de cada òrbita i determineu quantes voltes completes ha fet la nau al voltant del Sol en aquests 1\,000 dies. Quant mesura l'eix major de l'òrbita entre els moments A i C? (Doneu el resultat en radis solars.)
\end{apartat}

\begin{apartat}{1,25}
Representeu esquemàticament el Sol i l'òrbita de la nau entre els moments A i C. Indiqueu sobre el dibuix les posicions corresponents a A, B i C. Situeu la nau en la posició B i dibuixeu en aquest instant els vectors velocitat i acceleració de la nau (no cal calcular-ne els mòduls). En quina posició la velocitat de la nau és màxima? Justifiqueu la resposta i indiqueu el principi físic en què us baseu.
\end{apartat}
